\documentclass[10pt,a4paper]{amsart}
\usepackage{amsmath,amssymb,amsthm,graphicx}
\usepackage{verbatim}
\newenvironment{reference}{\comment}{\endcomment}
%\newenvironment{reference}{}{}
\newenvironment{slogan}{\comment}{\endcomment}
\newenvironment{history}{\comment}{\endcomment}

% For commutative diagrams we use Xy-pic
\usepackage[all]{xy}

% We use 2cell for 2-commutative diagrams.
\xyoption{2cell}
% Theorem environments.
%
\theoremstyle{plain}
\newtheorem{theorem}[subsection]{Theorem}
\newtheorem{proposition}[subsection]{Proposition}
\newtheorem{lemma}[subsection]{Lemma}

\theoremstyle{definition}
\newtheorem{definition}[subsection]{Definition}
\newtheorem{example}[subsection]{Example}
\newtheorem{exercise}[subsection]{Exercise}
\newtheorem{situation}[subsection]{Situation}

\theoremstyle{remark}
\newtheorem{remark}[subsection]{Remark}
\newtheorem{remarks}[subsection]{Remarks}

\numberwithin{equation}{subsection}

% Macros
%
\def\lim{\mathop{\mathrm{lim}}\nolimits}
\def\colim{\mathop{\mathrm{colim}}\nolimits}
\def\Spec{\mathop{\mathrm{Spec}}}
\def\Hom{\mathop{\mathrm{Hom}}\nolimits}
\def\Ext{\mathop{\mathrm{Ext}}\nolimits}
\def\SheafHom{\mathop{\mathcal{H}\!\mathit{om}}\nolimits}
\def\SheafExt{\mathop{\mathcal{E}\!\mathit{xt}}\nolimits}
\def\Sch{\mathit{Sch}}
\def\Mor{\mathop{\mathrm{Mor}}\nolimits}
\def\Ob{\mathop{\mathrm{Ob}}\nolimits}
\def\Sh{\mathop{\mathit{Sh}}\nolimits}
\def\NL{\mathop{N\!L}\nolimits}
\def\CH{\mathop{\mathrm{CH}}\nolimits}
\def\proetale{{pro\text{-}\acute{e}tale}}
\def\etale{{\acute{e}tale}}
\def\QCoh{\mathit{QCoh}}
\def\Ker{\mathop{\mathrm{Ker}}}
\def\Im{\mathop{\mathrm{Im}}}
\def\Coker{\mathop{\mathrm{Coker}}}
\def\Coim{\mathop{\mathrm{Coim}}}

% Boxtimes
%
\DeclareMathSymbol{\boxtimes}{\mathbin}{AMSa}{"02}

%
% Macros for moduli stacks/spaces
%
\def\QCohstack{\mathcal{QC}\!\mathit{oh}}
\def\Cohstack{\mathcal{C}\!\mathit{oh}}
\def\Spacesstack{\mathcal{S}\!\mathit{paces}}
\def\Quotfunctor{\mathrm{Quot}}
\def\Hilbfunctor{\mathrm{Hilb}}
\def\Curvesstack{\mathcal{C}\!\mathit{urves}}
\def\Polarizedstack{\mathcal{P}\!\mathit{olarized}}
\def\Complexesstack{\mathcal{C}\!\mathit{omplexes}}
% \Pic is the operator that assigns to X its picard group, usage \Pic(X)
% \Picardstack_{X/B} denotes the Picard stack of X over B
% \Picardfunctor_{X/B} denotes the Picard functor of X over B
\def\Pic{\mathop{\mathrm{Pic}}\nolimits}
\def\Picardstack{\mathcal{P}\!\mathit{ic}}
\def\Picardfunctor{\mathrm{Pic}}
\def\Deformationcategory{\mathcal{D}\!\mathit{ef}}
\title{A finite type map from an affine scheme to a quasi-separated algebraic space that admits no immersion into affine space}
\author{The Stacks Project Authors}
\date{stacks-project:examples.tex:local-snapshot-20261012; GNU FDL 1.2 or later}
\begin{document}
\maketitle
\noindent\textit{Editorial note.} All delivered mathematical text is the author's own, copied from the pinned Stacks Project chapter file examples.tex; no statement or proof step was rewritten, reordered or removed.\par\medskip
\noindent\textit{Editorial note.} The counted proof is the author's complete argument: The final 'which is not an immersion' is the author's own conclusion drawn from the displayed two-branch scheme, kept verbatim; the section's positive scheme statement (2402-2406) is provenance and is not part of the delivery.\par\medskip
\noindent\textit{Editorial note.} The label lemma-cannot-embed-into-affine names the counted statement, and the proof environment immediately following it in the source is delivered verbatim as its proof.\par\medskip
\noindent\textit{Editorial note.} The counted statement is carried by the author's own lemma environment in the chapter source.\par\medskip
\noindent\textit{Allowed standard prerequisites (editorial).}\begin{itemize}
\item Context span 2424-2436: Y = A\textasciicircum{}1\_k mapping to X = A\textasciicircum{}1\_k/R where R = \{(t, t)\} union \{(t, -t), t != 0\}; an immersion Y -> A\textasciicircum{}n\_X over X would be given by regular functions f\_1, ..., f\_n on Y; the fibre product of Y with the covering A\textasciicircum{}\{n+1\}\_k -> A\textasciicircum{}n\_X is the scheme \{(f\_1(t), ..., f\_n(t), t)\} union \{(f\_1(t), ..., f\_n(t), -t), t != 0\} with the obvious morphism to A\textasciicircum{}\{n+1\}\_k, which is not an immersion
\item The context closing sentence: this gives a counter example with X quasi-separated
\item The positive backdrop (section opening, 2402-2406): for schemes, a locally finite type morphism from an affine scheme immerses into some A\textasciicircum{}n\_X (Morphisms, quasi-affine-finite-type-over-S) - provenance only
\end{itemize}\medskip
A second counter example is $Y = \mathbf{A}^1_k \to X = \mathbf{A}^1_k/R$
with $R = \{(t, t)\} \amalg \{(t, -t), t \not = 0\}$. Namely, in
this case the morphism $Y \to \mathbf{A}^n_X$ would be given by some
regular functions $f_1, \ldots, f_n$ on $Y$ and hence the
fibre product of $Y$ with the covering
$\mathbf{A}^{n + 1}_k \to \mathbf{A}^n_X$
would be the scheme
$$
\{(f_1(t), \ldots, f_n(t), t)\} \amalg
\{(f_1(t), \ldots, f_n(t), -t), t \not = 0\}
$$
with obvious morphism to $\mathbf{A}^{n + 1}_k$ which is not an immersion.
Note that this gives a counter example with $X$ quasi-separated.

\begin{lemma}
\label{lemma-cannot-embed-into-affine}
There exists a finite type morphism of algebraic spaces $Y \to X$
with $Y$ affine and $X$ quasi-separated, such that there does not exist
an immersion $Y \to \mathbf{A}^n_X$ over $X$.
\end{lemma}

\begin{proof}
See discussion above.
\end{proof}

\end{document}
